%=============================================================================!
% 12.3  COUNTING STATES AT FIXED ENERGY                                       !
%=============================================================================!
\section{Counting states at fixed energy}
\label{sec:sm-micro}

A system that exchanges nothing with its surroundings keeps its energy,
as the runs of Chapters~9 to~11 did. This section gives the ensemble of
such a system, defines its entropy by counting its states, and finds that
temperature is what two systems sharing energy come to have in common.

\subsection*{The microcanonical ensemble}

Of all the clouds $\mathcal{P} = f(H)$ that the motion leaves unchanged, an
isolated system with energy $E$ calls for the one spread evenly over the
states of that energy, those with $H$ between $E$ and $E + \delta E$ for a
small $\delta E$, and zero elsewhere. If the system is ergodic, its
trajectory spends equal times near equal volumes of that shell, so the
even spread is the one its time averages sample. This is the
microcanonical ensemble. The volume of the shell, $\Omega(E)$, measured in
some fixed unit of phase-space volume, whose choice adds only a constant
to $S$,
counts the microstates of energy $E$ (it is not the bold matrix
$\boldsymbol\Omega$ of \cref{sec:in-reversal}), and Boltzmann's entropy is
\begin{equation}
   S = \kB\ln\Omega ,
   \label{eq:sm-entropy}
\end{equation}
with Boltzmann's constant $\kB = \SI{8.617333262e-5}{\electronvolt\per
\kelvin}$, the value SciPy gives. Taking the logarithm turns the counts of
two independent systems, which multiply, into entropies, which add.

\subsection*{Sharing energy}

Two systems placed in contact exchange energy, the total $E = E_1 + E_2$
staying fixed. Every microstate of the first combines with every
microstate of the second, so the pair has $\Omega_1(E_1)\,\Omega_2(E -
E_1)$ microstates in which the first holds $E_1$, and in the
microcanonical ensemble of the pair the probability of that division is
proportional to this number. The most probable division makes it largest,
or its logarithm, $\ln\Omega_1(E_1) + \ln\Omega_2(E - E_1)$, largest, which is
the same division since the logarithm rises whenever its argument does;
this needs its derivative with respect to $E_1$ to vanish:
\begin{equation*}
   \frac{\dd\ln\Omega_1}{\dd E_1} - \frac{\dd\ln\Omega_2}{\dd E_2} = 0 ,
\end{equation*}
by the chain rule, with $E_2 = E - E_1$ and $\dd E_2/\dd E_1 = -1$. The
two systems end where each has the same value of $\dd\ln\Omega/\dd E$,
which is therefore what they have in common in equilibrium. We call it
$1/\kB T$:
\begin{equation}
   \frac{1}{\kB T} = \frac{\dd\ln\Omega}{\dd E} , \qquad\text{or}\qquad
   \frac1T = \frac{\dd S}{\dd E} ,
   \label{eq:sm-temperature}
\end{equation}
which defines the temperature $T$, in kelvin. When the first system's
$\dd\ln\Omega/\dd E$ is the larger, moving energy into it raises the
product, so energy flows from the system with the smaller $1/T$, the
hotter one, to the colder, as it is found to.

\subsection*{The ideal gas}

$N$ atoms of mass $m$ in a box of volume $V$, with no forces between them,
have the energy $E = \sum_i\lvert\vb{p}_i\rvert^2/2m$ whatever their
positions. Each atom can be anywhere in the box, so the positions
contribute the factor $V^N$. The momenta with energy below $E$ are those
with $p_1^2 + p_2^2 + \dotsb + p_{3N}^2 \le 2mE$, listing all $3N$
components; by Pythagoras' theorem extended to $3N$ directions, these are
the points within the distance $R = \sqrt{2mE}$ of the origin, the inside
of a sphere of radius $R$ in $3N$ dimensions. That sphere is the sphere of
radius 1 with every coordinate stretched by the factor $R$, and
stretching all $3N$ coordinates by $R$ multiplies the volume of every
small box, the product of its $3N$ sides, by $R^{3N}$, so the volume
inside it is $cR^{3N}$, with $c$ the volume inside the sphere of radius 1,
a number that depends only on the dimension.
The number of states with energy below $E$ is therefore $cV^N(2mE)^{3N/2}$,
and those between $E$ and $E + \delta E$ number its derivative times
$\delta E$:
\begin{align*}
   \Omega(E) &= c\,V^N(2m)^{3N/2}\,\tfrac{3N}{2}E^{3N/2 - 1}\,\delta E ,\\
   \ln\Omega &= \bigl(\tfrac{3N}{2} - 1\bigr)\ln E + \text{a constant} ,\\
   \frac{1}{\kB T} &= \frac{3N/2 - 1}{E} \approx \frac{3N}{2E}
   && \text{by \cref{eq:sm-temperature}, for large $N$.}
\end{align*}
The third line uses the slope $1/E$ of $\ln E$: differentiating
$\mathrm{e}^{\ln E} = E$ by the chain rule gives $E\,\dd(\ln E)/\dd E =
1$.
The energy of the gas is $E = \frac32N\kB T$: $\frac32\kB T$ for each atom,
$\frac12\kB T$ for each direction of its motion, a sharing that
\cref{sec:sm-equipartition} finds to hold far more widely.

\subsection*{How sharp the sharing is}

Two such gases of $N$ atoms each, sharing $E$, have $\Omega_1\Omega_2
\propto E_1^{3N/2}(E - E_1)^{3N/2}$ for large $N$, so the probability that
the first holds the fraction $x = E_1/E$ of the energy is proportional to
$x^{3N/2}(1 - x)^{3N/2}$. It is largest at $x = \frac12$, equal energies
and equal temperatures, and \cref{fig:sm-sharing} shows how it narrows as
$N$ grows. Near the peak, the Taylor series of its logarithm
(\cref{sec:mo-taylor}) gives its width:
\begin{align*}
   \ln\mathcal{P} &= \tfrac{3N}{2}\bigl[\ln x + \ln(1 - x)\bigr]
   + \text{a constant} ,\\
   \frac{\dd\ln\mathcal{P}}{\dd x} &= \tfrac{3N}{2}\Bigl[\frac1x
   - \frac1{1 - x}\Bigr] = 0
   && \text{at $x = \tfrac12$,}\\
   \frac{\dd^2\ln\mathcal{P}}{\dd x^2} &= -\tfrac{3N}{2}\Bigl[\frac1{x^2}
   + \frac1{(1 - x)^2}\Bigr] = -12N
   && \text{at $x = \tfrac12$,}\\
   \ln\mathcal{P} &\approx \text{a constant} - \tfrac12\times12N
   \bigl(x - \tfrac12\bigr)^2 ,
\end{align*}
the chain rule giving the slope $-1/(1 - x)$ of $\ln(1 - x)$, and the
last line being the Taylor series to second order, whose first-order term
is zero at the peak. That last line is the logarithm of a Gaussian,
\cref{eq:sm-gauss}, with
$1/(2\sigma_x^2) = 6N$: the share has the standard deviation $\sigma_x =
1/\sqrt{12N}$. For 100 atoms in each gas that is 0.0289, against 0.0288
for the curve of \cref{fig:sm-sharing}, which counts the states in full,
computed numerically; for a thousand, 0.0091; for a macroscopic sample of
$10^{23}$ atoms, $\num{9e-13}$. A sample large enough to see is never
found far from its most probable division, which is why its temperature
is a definite number.

\begin{figure}[htb]
\centering
\includegraphics{ch12_ensembles/sharing}
\caption{Two gases of $N$ atoms each sharing a fixed energy: the
probability density for the first to hold the fraction $x$ of it, proportional to
$x^{3N/2 - 1}(1 - x)^{3N/2 - 1}$ by the full count $\Omega \propto
E^{3N/2 - 1}$, each curve divided by its largest value, for $N$ = 1, 10
and 100.}
\label{fig:sm-sharing}
\end{figure}

\begin{keyidea}
An isolated system, if it is ergodic, samples its energy shell evenly; the volume
$\Omega(E)$ of the shell counts its states, and $S = \kB\ln\Omega$ is its
entropy. Two systems sharing energy settle at the division with the most
states, where $\dd\ln\Omega/\dd E$, which defines $1/\kB T$, is the same
for both. An ideal gas holds $\frac32\kB T$ per atom, and the division of
energy between large systems is sharp, with a relative spread of order
$1/\sqrt N$.
\end{keyidea}

\notebook{12}{share energy between two gases of any size}

\begin{selfcheck}
For two gases of 10 atoms each, what does $1/\sqrt{12N}$ give for the
standard deviation of the share?
\end{selfcheck}
